
\documentclass{article} 
\usepackage{iclr2025_conference,times}


\input{math_commands.tex}

\usepackage{hyperref}
\usepackage{url}


\title{Formatting Instructions for ICLR 2025 \\ Conference Submissions}



\author{Raphaël Faure 2 \& Tristan Beruard \\
CentraleSupélec \\
\texttt{\{raphael.faure2, tristan.beruard\}@student-cs.fr} \\
}

\newcommand{\fix}{\marginpar{FIX}}
\newcommand{\new}{\marginpar{NEW}}

\iclrfinalcopy 
\begin{document}


\maketitle

\begin{abstract}

In our  

\end{abstract}

\section{Motivation and Problem Definition}

In 2026, \textit{cite} states that more than X\% of the code pushed to GitHub is AI generated. 
AI allows for faster proptype, and easier "launch of code" even in production. But \cite{} shows that without 
Careful code reviews of sometimes thousand of AI code lines, bugs and break can more frequently appears. 
Detecting whereas a code comes from an LLM or not can therefore reinforce the craefulness when reviewving it... 

Detecting AI generated content is a non trivial task and has been the subject of many research works. Many approaches use LLMs to detect AI generated content, and in particular, generated text.
Our ambition is not to detect AI generated content but to focus on a more specific content type : code.
There exist many different ways of representing code and it is certainly not limited to text representation. 
Very common representations include Abstract Syntax Tree (AST), Control Flow Graph (CFG) and Data Flow Graph (DFG). Those are widely
used in a broad range of applications such as code compilation, syntax error detection, and even vulnerability detection.

In this work, we will explore the use of graph-based representations of code, such as ASTs, CFGs, and DFGs, to detect AI-generated code.

\section{Methodology}


def of a System call 
graph formulation 

\section{Evaluation}

metric 

\section{References}


\subsubsection*{Author Contributions}
If you'd like to, you may include  a section for author contributions as is done
in many journals. This is optional and at the discretion of the authors.

\subsubsection*{Acknowledgments}
Use unnumbered third level headings for the acknowledgments. All
acknowledgments, including those to funding agencies, go at the end of the paper.


\bibliography{iclr2025_conference}
\bibliographystyle{iclr2025_conference}

\appendix
\section{Appendix}
https://arxiv.org/pdf/2509.16625
An Anomaly Intrusion Detection Method Based on PageRank Algorithm : https://ieeexplore-ieee-org.ezproxy.universite-paris-saclay.fr/document/6682431
dataset : https://old.dbs.uni-leipzig.de/file/Intrusion-Detection-on-System-Call-Graphs.pdf
more advanced dataset : https://old.dbs.uni-leipzig.de/file/BSI-LID-DS.pdf






\end{document}
